\documentclass[11pt]{article}
\usepackage{mathblatt}
\makeatletter
\setlength{\mbpfbe}{0mm}% keine Punkte (Bauregel 6.4)
\makeatother
% Rechenraum als Karo (Bauregel 6.9): n Rechenschritte = n mal 10 mm
\newcommand{\karo}[1]{\par\smallskip\noindent\begin{tikzpicture}
  \clip (0,0) rectangle ({\linewidth-1mm},{#1*10mm});
  \draw[black!16,line width=0.3pt,step=5mm] (0,0) grid ({\linewidth},{#1*10mm});
  \end{tikzpicture}\par}
\newcommand{\linie}[1]{\,{\color{black!45}\rule[-1pt]{#1}{0.4pt}}\,}
\newcommand{\pkt}{\textperiodcentered{}}
\begin{document}

\pfheftstil
\pflaufkopf[P10 \pkt{} mehr]{Netz und Schrägbild \pkt{} Lösungen}{PG80-4}\pfseitevon
\pfheftkopf{Netz und Schrägbild}{Lösungen \pkt{} mehr}
\begin{pfloesungb}{}
\lzb{1.}{}{B: Pyramide (zweite Option) – vier Dreiecke, drei davon klappen zur Spitze}{Pyramide (zweite Option) – vier Dreiecke, drei davon klappen zur Spitze}
\end{pfloesungb}
\begin{pfloesungb}{}
\lzb{2.}{}{Prisma}{Prisma}
\end{pfloesungb}
\begin{pfloesungb}{}
\lzb{3.}{}{h = 0,70 m; a = 0,90 m}{Höhe senkrecht von der Spitze zum Schnittpunkt der Diagonalen, gestrichelt; h = 0,70 m; a = 0,90 m}
\end{pfloesungb}
\begin{pfloesungb}{}
\lzb{4.}{}{Rückfläche 14 cm × 10 cm rechts an die Grundfläche; zweites Trapez unten an die Grundfläche}{Rückfläche 14 cm × 10 cm rechts an die Grundfläche; zweites Trapez unten an die Grundfläche}
\end{pfloesungb}
\begin{pfloesungb}{}
\lzb{5.}{}{Netz um Rechtecke 37 cm × 150 cm und 108 cm × 150 cm ergänzt; a = 150 cm}{Netz um Rechtecke 37 cm × 150 cm und 108 cm × 150 cm ergänzt; a = 108 cm, b = 37 cm, d = 150 cm}
\end{pfloesungb}
\end{document}